\chapter{The component library: Abstract}
\label{c:components}
\index{Library!components|textbf}

This chapter presents an abstract of the \MCS component library's structure
and categories. It intentionally does \emph{not} attempt to enumerate every
component with its parameter list -- the library is large and changes
between releases, so any hand-maintained table here would inevitably drift
out of date (as earlier editions of this manual did). Instead, each
component category chapter that follows (Sources, Optics, Samples,
Monitors, Misc, \ldots) documents a curated selection of representative
components in depth, with their full, current parameter lists pulled
directly from the component source's McDoc header at build time (see
section~\ref{s:mcdoc} for the McDoc comment format). For the complete,
authoritative, always-current list of every component and its parameters,
use the \verb+mcdoc+ front-end:
\begin{lstlisting}
mcdoc
mcdoc searchterm
mcdoc file.comp
\end{lstlisting}
The first form generates and opens a browsable \textit{index.html} catalog
of the whole library (both the standard installation and any local/custom
components); the second searches for and displays documentation for all
components matching \textit{searchterm}; the third displays the
documentation of one specific component file. See
section~\ref{s:mcdoc-run} in the User Manual for the full \verb+mcdoc+
reference.

\section{Component categories}
\label{s:comp-overview}

The library (located under the \verb+mcstas-comps+ directory of the \MCS
source/installation tree, see section~\ref{s:files} in the User Manual for
how \mcs locates it) is organised into the following top-level categories:

\begin{description}
\item[sources] Components that define the initial neutron state -- see
  chapter~\ref{c:source}.
\item[optics] Guides, slits, choppers, monochromators, mirrors, and other
  beam-manipulating components.
\item[samples] Components modelling matter placed in the beam for
  scattering studies (powders, single crystals, SANS models, \ldots) -- see
  chapter~\ref{c:samples}.
\item[monitors] Components that record/histogram the neutron beam without
  otherwise affecting it.
\item[misc] Components that do not fit the other categories, e.g. shape
  primitives, the \texttt{Progress\_bar}, and \texttt{MCPL\_input}/\texttt{MCPL\_output}
  for reading/writing portable neutron event files.
\item[union] The Union framework (contributed by Mads Bertelsen): a set of
  process/geometry components that can be freely combined to describe
  complex, possibly nested and overlapping sample environments with
  multiple scattering, decoupling the description of \emph{where} matter is
  from \emph{what} it does to the neutron. See the worked example in
  section~\ref{s:instrument} (Union\_demos) and the samples chapter.
\item[sasmodels] Small-angle scattering form factors generated from the
  SasView/sasmodels project, exposed via the \texttt{SasView\_model}
  component.
\item[contrib] Components contributed by \MCS users. The \MCS core team
  provides the same technical support as for other categories, but authored
  and maintained the code less directly; contributed components are
  generally reliable, but see each component's own validation/bugs notes.
  Components in \texttt{contrib} that mature and gain wide use are
  periodically promoted into one of the core categories above (e.g. several
  guide, chopper and polarisation components originally contributed have
  since moved into \texttt{optics}) -- always check with \verb+mcdoc+ rather
  than assuming a component's category from an older document.
\item[obsolete] Components that were renamed or superseded. They are kept,
  and continue to work, purely for backwards compatibility with existing
  instrument files; new instrument development should avoid them (\verb+mcdoc+
  will point to the current replacement where one exists).
\item[parked] Components that are kept in the source tree (e.g. pending
  further validation, a rewrite, or reinstatement) but excluded from the
  regular build/test cycle.
\end{description}

\section{Data files}
\index{Library!Components!data}

Some components require external data files, e.g. lattice crystallographic
definitions for Laue and powder pattern diffraction, $S(q,\omega)$ tables
for inelastic scattering, or transmission/reflectivity curves. Such files
distributed with \MCS are located in the \verb+data+ sub-directory of the
library and are accessible to any component using the \texttt{read\_table-lib}
shared library (see section~\ref{s:read-table}) without needing local
copies.

\begin{table}
  \begin{center}
    {\let\my=\\
    \begin{tabular}{|p{0.24\textwidth}|p{0.7\textwidth}|}
      \hline
       \textbf{MCSTAS/data} & Description \\
       \hline
 *.lau & Laue pattern file, as issued from Crystallographica.
       For use with Single\_crystal, PowderN, and Isotropic\_Sqw.
       Data: [ h   k   l Mult. d-space 2Theta   F-squared ] \\
 *.laz & Powder pattern file, as obtained from Lazy/ICSD.
       For use with PowderN, Isotropic\_Sqw and possibly Single\_crystal.\\
 *.trm & transmission file, typically for monochromator crystals and filters.
       Data: [ k (Angs-1) , Transmission (0-1) ] \\
 *.rfl & reflectivity file, typically for mirrors and monochromator crystals.
       Data: [ k (Angs-1) , Reflectivity (0-1) ] \\
 *.sqw & $S(q,\omega)$ files for Isotropic\_Sqw component.
       Data: [q] [$\omega$] [$S(q,\omega)$]\\
      \hline
    \end{tabular}
    \caption{Data files of the \MCS library.}
    \label{t:comp-data}
    \index{Library!Components!data}
    }
  \end{center}
\end{table}

\section{Component and instrument examples}

An overview of the example instrument library, organised by facility and
category, is given in the User Manual, chapter~\ref{s:instrument}. Every
\texttt{.instr} file there is discoverable via \verb+mcgui+'s
\textbf{New from Template} menu, or via \verb+mcdoc+.
